\documentclass[11pt]{article}

\usepackage[margin=1in]{geometry}
\usepackage{amsmath,amssymb}
\usepackage{array,longtable}
\usepackage[hidelinks]{hyperref}

\setlength{\parindent}{0pt}
\setlength{\parskip}{0.65em}
\renewcommand{\arraystretch}{1.18}

\title{A Lecture-Style Note on the BS20 Large-Economy Argument and the
Kline--Sorkin--Volpe Comparison}
\author{}
\date{}

\begin{document}
\maketitle

\section*{The two answers}

\textbf{BS20.}
The multiplier \(\tau\) is not something observed in the economy.  It
indexes a thought experiment in which we observe proportionally more
workers and firms generated by the same economic rules.  The number of jobs
per worker can stay small.  The estimator becomes accurate because it
averages noisy information across more distinct workers and firms.

\textbf{Kline, Sorkin, and Volpe.}
Kline's treatment-effect object
\(G_i(j,k)=Y_{i2}(k)-Y_{i2}(j)\) has the same mathematical form as the
worker-specific potential-wage contrast obtained by subtracting
Sorkin--Warwar's wage equation at firm \(j\) from that equation at firm
\(k\).  The assumptions must be compared separately.  Kline's use of
\emph{additive separability} matches
Volpe's wage-equation restriction: worker--firm interactions are absent.
Kline's no-selection-on-treatment-effects assumption is related to, but is
not mathematically equivalent to, Sorkin--Warwar's exogenous-mobility
assumption.  Kline's overall separation between a wage restriction and a
movement restriction matches Volpe's two-pillar organization.  Finally, an
average edge effect is a population average, not an assumption; it may have
an additive representation even when individual wages are not additively
separable.

\section*{How comparisons are made in this note}

To avoid comparing unlike things, the note keeps four categories separate:

\begin{enumerate}
\item An \textbf{object} is a wage, wage difference, average, or estimator.
\item An \textbf{assumption} restricts wages, movement, or sampling.
\item An \textbf{implication} is a conclusion proved from assumptions.
\item An \textbf{interpretation} explains the economic meaning of an object
or assumption.
\end{enumerate}

Objects are compared with objects, assumptions with assumptions, and
organizations of assumptions with organizations of assumptions.  An
interpretation is never treated as if it were a definition.

\section{What do \(\tau\) and \(I_x\) mean in BS20?}

\subsection{Notation}

\begin{longtable}{p{0.18\textwidth}p{0.54\textwidth}p{0.20\textwidth}}
\textbf{Symbol} & \textbf{Meaning in an actual labor-market dataset} &
\textbf{Source or status} \\
\hline
\(\mathcal X,\mathcal Y\) & Finite sets of possible worker and firm
characteristics in the statistical model. & BS20, p.\ 18; calligraphic
notation is a clarification. \\
\(x,y\) & A worker characteristic and a firm characteristic.  A
characteristic is the model's complete wage-, matching-, and
duration-relevant state; it need not be directly observed. & BS20,
pp.\ 18--20. \\
\(I_x,J_y\) & Numbers of distinct workers with characteristic \(x\) and
distinct firms with characteristic \(y\) in the starting theoretical
economy. & BS20, p.\ 18. \\
\(I,J\) & Total numbers of distinct workers and firms:
\(I=\sum_{x\in\mathcal X}I_x\) and
\(J=\sum_{y\in\mathcal Y}J_y\). & BS20, p.\ 18. \\
\(\tau\) & Positive-integer market-size multiplier used only in the
consistency proof. & BS20, p.\ 18. \\
\(i,j\) & A particular worker identity and firm identity. & BS20,
pp.\ 18--19. \\
\(x_i\) & Characteristic of worker \(i\). & BS20, pp.\ 18--21. \\
\(M_i,N_j\) & Numbers of distinct matches observed for worker \(i\) and firm
\(j\). & BS20, pp.\ 18--22. \\
\(m,m'\) & Two different matches in worker \(i\)'s history. & BS20,
pp.\ 19--22. \\
\(w^w_{i,m}\) & Average log wage in worker \(i\)'s match \(m\). & BS20,
p.\ 19. \\
\(\lambda(x_i)\) & Expected log wage received by a worker with
characteristic \(x_i\); this is the worker's population wage type. & BS20,
pp.\ 5, 20--22. \\
\(\mu(y)\) & Expected log wage paid by a firm with characteristic \(y\);
this is the firm's population wage type. & BS20, pp.\ 5, 20--23. \\
\(\varepsilon^w_{i,m}\) & Match-level deviation of
\(w^w_{i,m}\) from \(\lambda(x_i)\) in the worker-side auxiliary equation. &
BS20, pp.\ 21--22. \\
\(S_i\) & Noisy worker-level score whose expectation is
\(\lambda(x_i)^2\). & Derived shorthand for BS20 equation (18), p.\ 22. \\
\(\eta_i\) & Score error \(S_i-\lambda(x_i)^2\). & Derived notation. \\
\(A_\tau\) & Average score error across the \(\tau I\) workers. & Derived
notation. \\
\(C,\delta\) & Finite variance bound and arbitrary positive tolerance in
the probability proof.  These are proof devices, not economic variables. &
Standard proof notation. \\
\(\rho_{\mathrm{BS}},\widehat\rho_{\mathrm{BS}}\) & Population sorting
correlation between matched worker and firm expected-wage types, and its
sample estimator. & BS20, pp.\ 5--6, 24. \\
\end{longtable}

\subsection{Start with the actual objects}

The law of large numbers (LLN) used below has one simple role: when many
independent errors each have mean zero and finite variance, their average
converges to zero.  BS20 constructs unit-level errors with mean zero and
then increases the number of units being averaged.

Suppose there are several possible worker characteristics.  The starting
theoretical economy contains \(I_x\) workers of each characteristic \(x\).
It therefore contains
\[
I=\sum_{x\in\mathcal X}I_x
\tag{B.1}
\]
workers in total.  Similarly, it contains \(J_y\) firms of each firm
characteristic \(y\), for a total of
\[
J=\sum_{y\in\mathcal Y}J_y.
\tag{B.2}
\]

For example, a starting economy might contain 60 workers in one
characteristic group and 40 workers in another, so \(I=100\).  The workers
within one group share the same expected-wage type, but they are different
people.  They can experience different employers, wages, and match
durations.

The phrase \emph{starting economy} means only the finite template from which
BS20 constructs its mathematical sequence.  It need not be a literal
country, and it is not an earlier time period.

\textit{Source: BS20, p.\ 18.  The numerical example and the term
``finite template'' are explanatory interpretations.}

\subsection{What the multiplier \(\tau\) does}

BS20 constructs a larger hypothetical economy by multiplying every
characteristic-specific population count by the same number \(\tau\).  The
larger economy contains
\[
\tau I_x
\quad\text{workers of characteristic }x
\tag{B.3}
\]
and
\[
\tau J_y
\quad\text{firms of characteristic }y.
\tag{B.4}
\]
Its total numbers of workers and firms are therefore
\[
\tau I
\qquad\text{and}\qquad
\tau J.
\tag{B.5}
\]

The population shares do not change:
\[
\frac{\tau I_x}{\tau I}=\frac{I_x}{I},
\qquad
\frac{\tau J_y}{\tau J}=\frac{J_y}{J}.
\tag{B.6}
\]
Thus \(\tau\) makes the market larger without changing its composition.

If the starting economy has 60 workers in one group and 40 in another, then
\(\tau=10\) produces 600 workers in the first group and 400 in the second.
It does not produce ten observations on each original worker.  It produces
new worker identities whose employment histories are new draws from the
same group-specific statistical rules.

This is why BS20 uses a multiplier instead of merely writing
``\(I\to\infty\) and \(J\to\infty\).''  The multiplier specifies exactly
how both sides of the market grow and guarantees that the fractions of each
worker and firm characteristic remain fixed.

\textit{Source: BS20, p.\ 18.  Equations (B.3)--(B.6) expand the paper's
replication statement.}

\subsection{Workers with the same characteristic are still different draws}

For a worker \(i\) with characteristic \(x_i\), BS20's auxiliary wage
equation is
\[
w^w_{i,m}
=
\lambda(x_i)+\varepsilon^w_{i,m},
\qquad
\mathbb E[\varepsilon^w_{i,m}\mid x_i]=0.
\tag{B.7}
\]
Two workers with the same \(x_i\) share the same expected component
\(\lambda(x_i)\), but their errors and match histories differ.  Having the
same characteristic therefore means having the same conditional
distribution, not having the same realized data.

This is analogous to drawing repeatedly from the same distribution.  Every
draw follows the same rule, but the realized outcomes differ.  More
independent draws reveal the common moments of that distribution more
accurately.

\textit{Source: BS20, pp.\ 20--22 and Auxiliary Assumption 1.  The
same-distribution interpretation is explanatory.}

\subsection{The LLN proof}

\textbf{Within-worker assumptions.}
Assume worker \(i\) has at least two qualifying matches.  For different
matches \(m\neq m'\),
\[
\begin{aligned}
&\mathbb E[
w^w_{i,m}w^w_{i,m'}\mid x_i]\\
&=
\mathbb E[
(\lambda(x_i)+\varepsilon^w_{i,m})
(\lambda(x_i)+\varepsilon^w_{i,m'})
\mid x_i]\\
&=
\lambda(x_i)^2
+\lambda(x_i)
\mathbb E[
\varepsilon^w_{i,m}+\varepsilon^w_{i,m'}\mid x_i]\\
&\quad+
\mathbb E[
\varepsilon^w_{i,m}\varepsilon^w_{i,m'}\mid x_i]\\
&=
\lambda(x_i)^2.
\end{aligned}
\tag{B.8}
\]
The two single-error terms are zero by conditional mean zero.  Conditional
independence of different-match errors makes the error-product expectation
zero.

\textbf{Worker-level object.}
Average all ordered products of different matches:
\[
S_i
=
\frac{1}{M_i(M_i-1)}
\sum_{m=1}^{M_i}
\sum_{\substack{m'=1\\m'\neq m}}^{M_i}
w^w_{i,m}w^w_{i,m'}.
\tag{B.9}
\]
Equation (B.8) gives
\[
\mathbb E[S_i\mid x_i]=\lambda(x_i)^2.
\tag{B.10}
\]
The score \(S_i\) is unbiased, but it remains noisy if \(M_i\) remains
small.  BS20 does not claim that each worker's type becomes known exactly.

\textbf{Across-worker assumptions.}
Define the worker-level score error and its cross-worker average:
\[
\eta_i=S_i-\lambda(x_i)^2,
\qquad
A_\tau=\frac{1}{\tau I}\sum_{i=1}^{\tau I}\eta_i.
\tag{B.11}
\]
Under the cross-worker independence and finite-variance conditions used in
the BS20 proof,
\[
\mathbb E[\eta_i\mid x_i]=0,
\qquad
\operatorname{Var}(\eta_i)\leq C<\infty.
\tag{B.12}
\]

\textbf{LLN implication.}
Therefore,
\[
\begin{aligned}
\operatorname{Var}(A_\tau)
&=
\frac{1}{(\tau I)^2}
\sum_{i=1}^{\tau I}\operatorname{Var}(\eta_i)\\
&\leq
\frac{\tau I C}{(\tau I)^2}\\
&=
\frac{C}{\tau I}
\longrightarrow0.
\end{aligned}
\tag{B.13}
\]
Chebyshev's inequality then gives, for every \(\delta>0\),
\[
\Pr(|A_\tau|>\delta)
\leq
\frac{C}{\tau I\delta^2}
\longrightarrow0.
\tag{B.14}
\]
Thus the average score error disappears as the number \(\tau I\) of
distinct workers grows.  The error need not disappear for any particular
worker.

\textbf{Extension to the full estimator.}
The firm-side proof is the same, averaging noisy firm scores across
\(\tau J\) firms.  BS20's covariance estimator uses leave-one-match-out
worker and firm wage averages.  The paper first shows that their product has
the correct expectation and then averages those products across the growing
economy.  These consistent component estimates are combined to obtain
\[
\widehat\rho_{\mathrm{BS}}
\overset p\longrightarrow
\rho_{\mathrm{BS}}.
\tag{B.15}
\]
Positive worker- and firm-type variances are required for the correlation to
be defined.

In the actual dataset BS20 computes the formulas at \(\tau=1\).
The sequence \(\tau\to\infty\) exists only in the proof.  It says that one
large economy provides many distinct worker- and firm-level scores, even
when each individual has only a few matches.

\textit{Source: BS20, pp.\ 21--24, Propositions 3--5; Appendix B.1,
pp.\ 48--50.  Equations (B.8)--(B.14) isolate the LLN in elementary form.}

\subsection{What if there is only one worker characteristic?}

The LLN can still estimate the common worker moment because the economy
contains many different workers with independent employment histories.
However, if every worker has the same characteristic, every worker has the
same expected-wage type.  The worker-type variance is then zero.  The BS20
sorting correlation is not defined because a correlation requires positive
variation on both sides of the market.

Thus two statements must be kept separate:

\begin{enumerate}
\item Many workers with the same characteristic are useful for estimating
the moment associated with that characteristic.
\item Several characteristics with different expected-wage types are needed
for a nondegenerate sorting correlation.
\end{enumerate}

\textit{Derived from the definition of a correlation and the BS20
replication scheme.}

\section{Which definition is Kline closer to: Sorkin's or Volpe's?}

\subsection{Notation}

\begin{longtable}{p{0.18\textwidth}p{0.54\textwidth}p{0.20\textwidth}}
\textbf{Symbol} & \textbf{Meaning} & \textbf{Source or status} \\
\hline
\(i,j,k,t,s,T\) & Worker, origin firm, destination firm, outcome period,
conditioning period, and total number of periods. & Standard indices used
by the three sources. \\
\(D_{it}\) & Firm employing worker \(i\) in period \(t\). & Kline,
pp.\ 29--33. \\
\(\mathbf D_i\) & Worker \(i\)'s two-period firm path
\((D_{i1},D_{i2})\). & Derived shorthand. \\
\(Y_{it}(j),Y_{it}\) & Potential log wage at firm \(j\), and observed log
wage \(Y_{it}=Y_{it}(D_{it})\). & Kline, pp.\ 29--33. \\
\(X_{it},\beta\) & Worker-time covariates and their coefficient. &
Sorkin--Warwar, p.\ 5; Volpe, slide 4/35. \\
\(\alpha_i,\psi_j\) & Worker component and firm component of wages. &
Sorkin--Warwar, pp.\ 5--7; Volpe, slide 4/35; Kline, p.\ 15. \\
\(\varpi_{ij}\) & Permanent worker--firm wage interaction.  This relabels
Sorkin--Warwar's \(\mu_{ij}\) to avoid overloading \(\mu\). &
Sorkin--Warwar, p.\ 6; symbol relabeled. \\
\(\widetilde\varepsilon_{it}(j)\) & Transitory potential-wage remainder at
firm \(j\). & Sorkin--Warwar, pp.\ 6--7; firm index made explicit. \\
\(\varepsilon^A_{it}\) & Residual in the observed additive AKM equation. &
Derived from Sorkin--Warwar equation (2), p.\ 6. \\
\(d,\mathbf x\) & Generic firm and covariate values used when conditioning
on movement. & Derived notation. \\
\(\mathsf A_0\) & Pure additive-separability restriction
\(\varpi_{ij}=0\) for every \(i,j\). & Derived label. \\
\(\mathsf{EM}_{SW},\mathsf{AS}_{SW}\) & Sorkin--Warwar's named exogenous
mobility and additive-separability assumptions. & Derived labels for
Sorkin--Warwar, pp.\ 6--7. \\
\(\mathsf{EM}_{V},\mathsf{AS}_{V}\) & Volpe's exogenous-mobility and
additive-separability pillars. & Derived labels for Volpe, slide 4/35. \\
\(G_i(j,k)\) & Worker \(i\)'s wage treatment effect from replacing firm
\(j\) with firm \(k\). & Derived shorthand for Kline, pp.\ 30--33. \\
\(\Delta_{jk}\) & Average treatment effect for workers who move on the
directed edge \(j\to k\). & Kline, pp.\ 30--31. \\
\(\psi_j^{\mathrm{ATE}}\) & Population average potential-wage level at firm
\(j\), relative to reference firm \(0\). & Derived relabeling of Kline,
p.\ 33. \\
\(\xi_i\) & Nondegenerate, mean-zero worker-specific match heterogeneity in the
counterexample. & Derived notation. \\
\(\chi_i\) & Worker-specific period-2 wage shock that enters every potential
firm wage equally in the first non-equivalence example. & Derived notation. \\
\end{longtable}

\subsection{The wage equation and two different kinds of assumptions}

\textbf{Wage object.}
Write worker \(i\)'s potential wage at firm \(j\) as
\[
Y_{it}(j)
=
\alpha_i+\psi_j+X_{it}'\beta
+\varpi_{ij}
+\widetilde\varepsilon_{it}(j).
\tag{K.1}
\]
The permanent interaction \(\varpi_{ij}\) allows the value of firm \(j\) to
depend on which worker is employed there.  If an additive AKM regression
omits this interaction, its observed residual is
\[
\varepsilon^A_{it}
=
\varpi_{iD_{it}}
+\widetilde\varepsilon_{it}(D_{it}).
\tag{K.2}
\]

\textbf{Two kinds of assumptions.}
There are now two different questions:

\begin{enumerate}
\item \textbf{Wage question:} Is the interaction
\(\varpi_{ij}\) absent?
\item \textbf{Movement question:} Are workers' firm choices related to the
unobserved wage components?
\end{enumerate}

The first question concerns how wages are produced.  The second concerns who
moves where.  An assumption answering one question does not automatically
answer the other.

\textit{Source: Sorkin--Warwar, pp.\ 5--7, equations (1)--(2).
Equation (K.1) uses potential-wage notation and relabels their match effect.}

\subsection{Sorkin--Warwar's two definitions}

\textbf{Definitions.}
For every relevant worker, firm, period, and covariate cell, write the two
mean-exogeneity conditions as
\[
\mathbb E[
\widetilde\varepsilon_{it}(j)
\mid D_{is}=d,X_{is}=\mathbf x]
=0
\tag{K.3}
\]
and
\[
\mathbb E[
\varpi_{ij}
\mid D_{is}=d,X_{is}=\mathbf x]
=0.
\tag{K.4}
\]
Sorkin--Warwar call the combination of (K.3) and (K.4) exogenous mobility:
\[
\mathsf{EM}_{SW}
=
\text{(K.3) and (K.4)}.
\tag{K.5}
\]
They call the combination of (K.3) and
\[
\mathsf A_0:
\quad
\varpi_{ij}=0
\quad\text{for every }i,j
\tag{K.6}
\]
additive separability:
\[
\mathsf{AS}_{SW}
=
\text{(K.3) and (K.6)}.
\tag{K.7}
\]

\textbf{Implication.}
If \(\varpi_{ij}=0\), its conditional mean is automatically zero.  Therefore
\[
\mathsf{AS}_{SW}
\quad\Longrightarrow\quad
\mathsf{EM}_{SW}.
\tag{K.8}
\]
The nesting occurs because Sorkin's named additive-separability assumption
already includes the transitory-error condition (K.3).

\textbf{Interpretation.}
Sorkin's lecture summarizes the economic distinction as:

\begin{center}
exogenous mobility permits average effects; additive separability imposes
constant effects.
\end{center}

\textit{Sources: Sorkin--Warwar, pp.\ 6--7, Assumptions 1--2; Sorkin
Lecture 10, slide 46/79.  The implication in (K.8) is derived.}

\subsection{Volpe's two definitions}

\textbf{Additive-separability definition.}
Volpe uses \emph{additive separability} for the wage restriction alone:
\[
\mathsf{AS}_{V}:
\quad
\varpi_{ij}=0
\quad\text{for every }i,j.
\tag{K.9}
\]

\textbf{Exogenous-mobility definition.}
Volpe separately uses \emph{exogenous mobility} for the requirement that
movement not be driven by unobserved wage fluctuations.  The slide's panel
condition can be written
\[
\mathsf{EM}_{V}:
\quad
\mathbb E\left[
\varepsilon^A_{it}
\mid
\alpha_i,\{\psi_{D_{is}},X_{is}\}_{s=1}^{T}
\right]
=0
\quad\text{for every }t.
\tag{K.10}
\]

Condition (K.9) restricts the wage equation.  Condition (K.10) restricts
movement relative to the residual.  Neither condition automatically gives
the other.  This is why Volpe describes them as two pillars required for
the usual AKM interpretation.

\textit{Source: Volpe, slide 4/35.  The symbols
\(\mathsf{AS}_{V}\) and \(\mathsf{EM}_{V}\) are derived labels.}

\subsection{Kline's wage-difference object and selection assumption}

\textbf{Object.}
Kline considers the same worker's potential wage at two firms.  Define
\[
G_i(j,k)
=
Y_{i2}(k)-Y_{i2}(j).
\tag{K.11}
\]
This is a treatment effect: it is the wage change worker \(i\) would receive
if firm \(j\) were replaced by firm \(k\).  Only one of the two potential
wages is normally observed, but both are well-defined in the causal model.

\textbf{Derived identity.}
Substituting (K.1) into (K.11) gives
\[
\begin{aligned}
G_i(j,k)
&=
\psi_k-\psi_j\\
&\quad+
\varpi_{ik}-\varpi_{ij}\\
&\quad+
\widetilde\varepsilon_{i2}(k)
-\widetilde\varepsilon_{i2}(j).
\end{aligned}
\tag{K.12}
\]
The worker component and worker-time covariates cancel because the same
worker at the same time is compared at two firms.

Under additive separability, \(\varpi_{ij}=0\).  Removing the transitory
remainder from (K.12), every worker then has the same permanent wage gain:
\[
G_i(j,k)
-\left[
\widetilde\varepsilon_{i2}(k)
-\widetilde\varepsilon_{i2}(j)
\right]
=
\psi_k-\psi_j.
\tag{K.13}
\]
This is the constant-effect interpretation.

\textbf{Assumption.}
Kline's no-selection-on-treatment-effects assumption instead says
\[
G_i(j,k)\perp\mathbf D_i,
\qquad
\mathbf D_i=(D_{i1},D_{i2}).
\tag{K.14}
\]
In words, workers who move from \(j\) to \(k\) must not have systematically
different gains from that move than the population.  Gains may still vary
across workers.  Kline restricts who selects onto an edge; he does not set
the worker--firm interaction to zero.

\textit{Source: Kline, pp.\ 32--33, Assumption 4.  Equations
(K.11)--(K.13) are derived from the common wage decomposition.}

\subsection{Three comparisons that must be kept separate}

\textbf{Comparison 1: object with object.}
Kline's object is the worker-specific potential-wage difference
\[
G_i(j,k)=Y_{i2}(k)-Y_{i2}(j).
\]
Sorkin--Warwar analyze the same kind of object when they distinguish a
worker-specific firm effect from a constant firm effect: subtracting their
potential-wage equation at \(j\) from the equation at \(k\) produces the
same worker-specific contrast.  Volpe's slide begins instead from a
wage-level equation; a two-firm wage difference can be derived from it, but
it is not the slide's organizing object.

\textbf{Conclusion for Comparison 1.}
Kline's object and the contrast derived from Sorkin--Warwar's wage equation
have the same mathematical form.  This is an object-to-object statement; it
is not a claim that their assumptions are identical.

\textbf{Comparison 2: additive-separability definition with
additive-separability definition.}
Kline uses additive separability to mean that unobserved worker and firm
heterogeneity enter wages additively.  In (K.1), that is the wage restriction
\[
\varpi_{ij}=0
\quad\text{for every }i,j.
\]
This is also Volpe's definition in (K.9).  Sorkin--Warwar's named
additive-separability assumption contains this wage restriction
\emph{and} the transitory-error condition (K.3).

\textbf{Conclusion for Comparison 2.}
Kline's additive-separability definition matches Volpe's pure wage
restriction.  It differs from Sorkin--Warwar's named assumption because the
latter also contains the transitory-error condition.

\textbf{Comparison 3: organization of assumptions with organization of
assumptions.}
Kline treats the zero-interaction wage restriction and the no-selection
condition (K.14) as different possible restrictions.  Volpe likewise
presents additive separability and exogenous mobility as two separate
pillars.  Sorkin--Warwar instead define two packages that share (K.3), which
is why their named assumptions are nested.

\textbf{Conclusion for Comparison 3.}
Kline's organization matches Volpe's two-pillar organization.  This does not
make Kline's no-selection condition equal to Volpe's exogenous-mobility
condition; it says only that both authors separate a wage restriction from a
movement restriction.

\begin{longtable}{p{0.27\textwidth}p{0.28\textwidth}p{0.36\textwidth}}
\textbf{Level of comparison} & \textbf{Correct comparison} &
\textbf{Exact reason} \\
\hline
Potential-wage-difference object & Kline and Sorkin--Warwar & Both compare
the same worker's potential wage across two firms. \\
Definition of additive separability & Kline and Volpe & Both use the pure
wage restriction that rules out the permanent worker--firm interaction. \\
Organization of assumptions & Kline and Volpe & Both keep the wage
restriction separate from the movement restriction. \\
Movement assumption & No exact match & Kline restricts the full
distribution of a wage difference; Sorkin--Warwar restrict conditional means
of wage-level components; Volpe restricts the composite AKM residual. \\
\end{longtable}

\textit{Sources: Kline, pp.\ 30--33; Sorkin--Warwar, pp.\ 6--7;
Sorkin Lecture 10, slide 46/79; Volpe, slide 4/35.  The three comparison
statements are derived mappings.}

\subsection{Kline's and Sorkin--Warwar's movement assumptions are related,
not equivalent}

The two assumptions address a similar economic concern: workers who select
particular firms may have unusual unobserved wage gains.  Their mathematical
statements nevertheless differ.

\textbf{Kline's assumption.}
For each fixed pair \(j,k\), the entire distribution of the difference
\(G_i(j,k)\) is independent of the two-period mobility path:
\[
G_i(j,k)\perp\mathbf D_i.
\tag{K.14a}
\]

\textbf{Sorkin--Warwar's assumption.}
The conditional means of the level components
\(\widetilde\varepsilon_{it}(j)\) and \(\varpi_{ij}\) are zero in the
specified movement and covariate cells, as stated in (K.3)--(K.4).

\textbf{Why Kline's condition need not imply Sorkin--Warwar's.}
Set the permanent interactions to zero and suppose the period-2 transitory
remainder is common across potential firms:
\[
\widetilde\varepsilon_{i2}(j)=\chi_i
\quad\text{for every firm }j,
\tag{K.14b}
\]
and \(\chi_i\) is related to the worker's mobility path.  The common
component cancels from every two-firm difference:
\[
\widetilde\varepsilon_{i2}(k)
-\widetilde\varepsilon_{i2}(j)
=
\chi_i-\chi_i
=0.
\tag{K.14c}
\]
Kline's gain is therefore the constant \(\psi_k-\psi_j\), so Kline's
independence condition holds exactly.  But
\(\mathbb E[\widetilde\varepsilon_{i2}(j)\mid\mathbf D_i]
=\mathbb E[\chi_i\mid\mathbf D_i]\)
can be nonzero and vary with the mobility path, violating
Sorkin--Warwar's transitory-error condition.

\textbf{Why Sorkin--Warwar's mean condition need not imply Kline's full
independence.}
Set the transitory remainders to zero.  For two firms \(j,k\), set
\(\varpi_{ij}=0\) in two mobility groups.  Let
\(\varpi_{ik}\) take values \(-1\) and \(1\), each with probability one
half, in the first group, and values \(-2\) and \(2\), each with probability
one half, in the second group.  In each group,
\[
\mathbb E[\varpi_{ij}\mid\mathbf D_i]
=
\mathbb E[\varpi_{ik}\mid\mathbf D_i]
=0.
\tag{K.14d}
\]
Thus Sorkin--Warwar's level mean restriction holds.  But the distribution
of \(\varpi_{ik}-\varpi_{ij}\), and hence the distribution of the wage gain,
depends on the mobility group.  Kline's independence condition fails.  A
zero conditional mean therefore does not imply independence.

\textbf{Conclusion.}
Kline and Sorkin--Warwar restrict different mathematical objects and use
different strengths of restriction.  It is accurate to say that they
address related selection concerns.  It is inaccurate to call the
assumptions equivalent.

\textit{Source for the assumptions: Kline, p.\ 32, Assumption 4;
Sorkin--Warwar, p.\ 6, Assumption 1.  Equations (K.14b)--(K.14c) and both
non-equivalence examples are derived.}

\subsection{The average edge effect is an object, not an assumption}

\textbf{Object.}
Kline's average edge effect is
\[
\Delta_{jk}
=
\mathbb E[
G_i(j,k)
\mid D_{i1}=j,D_{i2}=k].
\tag{K.15}
\]
It is the average wage gain for the workers who actually move from \(j\) to
\(k\).

\textbf{Identification result.}
Under Kline's exclusion, parallel-trends, and stationarity assumptions, the
mean observed wage change on the edge equals the object in (K.15).  Those
three assumptions give the edge average a causal meaning.  They do not
require additive separability.

\textbf{Implication of Kline's no-selection assumption.}
If the no-selection condition (K.14) also holds, then
\[
\begin{aligned}
\Delta_{jk}
&=
\mathbb E[
G_i(j,k)
\mid D_{i1}=j,D_{i2}=k]\\
&=
\mathbb E[G_i(j,k)]\\
&=
\mathbb E[Y_{i2}(k)]
-\mathbb E[Y_{i2}(j)].
\end{aligned}
\tag{K.16}
\]
Define the average wage level at firm \(j\), relative to firm \(0\), as
\[
\psi_j^{\mathrm{ATE}}
=
\mathbb E[Y_{i2}(j)]
-\mathbb E[Y_{i2}(0)].
\tag{K.17}
\]
Then
\[
\Delta_{jk}
=
\psi_k^{\mathrm{ATE}}
-\psi_j^{\mathrm{ATE}}.
\tag{K.18}
\]

\textbf{Why this does not define additive separability.}
Two different assumptions can produce an additive-looking edge average:

\begin{enumerate}
\item Under Sorkin--Warwar's additive-separability package, the permanent
gain is \(\psi_k-\psi_j\) for every worker and the transitory remainder has
zero conditional mean.  Averaging therefore gives
\(\Delta_{jk}=\psi_k-\psi_j\).
\item Under Kline's no-selection assumption, individual gains may differ,
but the movers on every edge have the population gain distribution.
Averaging therefore gives
\(\Delta_{jk}
=\psi_k^{\mathrm{ATE}}-\psi_j^{\mathrm{ATE}}\).
\end{enumerate}

The same form for an average can therefore arise from constant individual
gains or from heterogeneous individual gains with no selection on those
gains.  Equation (K.18) is an implication of Kline's selection assumption,
not a definition of additive separability and not evidence that every
individual wage is additive.

\textit{Source: Kline, pp.\ 30--33, Propositions 1--2.  Equations
(K.15)--(K.18) write the argument line by line.}

\subsection{A short proof that average additivity is not individual
additivity}

\textbf{Construction.}
Consider two firms, \(0\) and \(1\), and set covariates and transitory
errors to zero.  Let
\[
\varpi_{i0}=\xi_i,
\qquad
\varpi_{i1}=-\xi_i,
\qquad
\mathbb E[\xi_i]=0,
\tag{K.19}
\]
where \(\xi_i\) is nondegenerate and independent of the mobility path
\(\mathbf D_i\).  Potential wages are
\[
Y_{i2}(0)=\alpha_i+\psi_0+\xi_i,
\qquad
Y_{i2}(1)=\alpha_i+\psi_1-\xi_i.
\tag{K.20}
\]
Therefore,
\[
G_i(0,1)
=
\psi_1-\psi_0-2\xi_i.
\tag{K.21}
\]

\textbf{Implication for individual wages.}
Because \(\xi_i\) is nondegenerate, individual treatment effects are
heterogeneous.  Moreover, the interaction terms in (K.19) are not
identically zero, so additive separability fails.

\textbf{Implication for selection and the edge average.}
Because \(\xi_i\perp\mathbf D_i\), equation (K.21) implies Kline's
no-selection condition.  Taking the mean gives
\[
\begin{aligned}
\Delta_{01}
&=
\mathbb E[G_i(0,1)]\\
&=
\psi_1-\psi_0-2\mathbb E[\xi_i]\\
&=
\psi_1-\psi_0.
\end{aligned}
\tag{K.22}
\]

\textbf{Conclusion.}
The average edge effect is an additive firm difference even though the
individual wage equation contains a nonzero worker--firm interaction.
Therefore,
\[
\boxed{
\text{additive average edge effects do not imply individual additive
separability}.
}
\tag{K.23}
\]

\textit{Equations (K.19)--(K.23) are a derived counterexample.}

\section*{References and page convention}

\begin{enumerate}
\item Borovi\v{c}kov\'a, Katar\'ina, and Robert Shimer.
\emph{High Wage Workers Work for High Wage Firms}, February 13, 2020.
Page references use the paper's printed pages.
\item Kline, Patrick M.  \emph{Firm Wage Effects}.  NBER Working Paper
33084, October 2024, revised February 2025.  Page references use the
paper's printed pages.
\item Sorkin, Isaac, and Lucas Warwar.  \emph{Paired Movers}, March 2026.
Page references use the paper's printed pages.
\item Sorkin, Isaac.  \emph{Labor Economics, Lecture 10}.  The cited slide
is numbered 46/79.
\item Volpe, Oscar.  \emph{Labor Omnibus Lectures, November 3 and 5}.
The cited slide is numbered 4/35.
\end{enumerate}

\end{document}
